\documentclass[11pt,a4paper]{article}
\usepackage[T1]{fontenc}
\usepackage[utf8]{inputenc}
\usepackage{amsmath,amssymb,amsthm}
\usepackage[margin=1in]{geometry}
\usepackage[colorlinks=true,linkcolor=blue,urlcolor=blue]{hyperref}
\theoremstyle{plain}
\newtheorem{theorem}{Theorem}
\newtheorem{lemma}[theorem]{Lemma}
\newtheorem{proposition}[theorem]{Proposition}
\newtheorem{corollary}[theorem]{Corollary}
\theoremstyle{definition}
\newtheorem{remark}[theorem]{Remark}

\title{The empirical recurrence for OEIS A250941, proved by transfer matrix}
\author{Adrian Perez Fontelles\\ \small Independent researcher}
\date{3 September 2026}
\begin{document}
\maketitle

\begin{abstract}
OEIS A250941 counts arrays over $\{0,1,2\}$ in which a statistic taken over short runs
of cells produces a derived array that is monotone in a named direction. The entry carries an
empirical recurrence of order $39$ contributed by R.~H.~Hardin. It is true, and it is
decidable rather than empirical: every comparison the condition makes lies inside a bounded
window of consecutive rows, so the count is a walk count on $S=6561$ states, and the
conjectured recurrence is then settled by a finite exact computation.
\end{abstract}

\noindent\small 2020 Mathematics Subject Classification. 05A15, 15A18.\normalsize

\section{The conjecture}

OEIS A250941 is ``Number of (n+1)X(7+1) 0..2 arrays with nondecreasing x(i,j)+x(i,j-1) in the i direction and nondecreasing absolute value of x(i,j)-x(i-1,j) in the j direction.''. It has offset $1$ and begins
\[
103093, 1188599, 9865706, 72722706, 475951361, 2839711239, 15686582837, 81197881457,\ \dots
\]
The entry states, as a conjecture contributed by its author and never marked settled:
\begin{quote}
Empirical: a(n) = 56*a(n-1) -1522*a(n-2) +26750*a(n-3) -341788*a(n-4) +3383870*a(n-5) -27019370*a(n-6) +178816330*a(n-7) -1000398245*a(n-8) +4801705370*a(n-9) -19999827232*a(n-10) +72938367572*a(n-11) -234575011664*a(n-12) +669091458700*a(n-13) -1700431069076*a(n-14) +3864487952276*a(n-15) -7876529587291*a(n-16) +14429295594092*a(n-17) -23796815540850*a(n-18) +35368238093958*a(n-19) -47398136192484*a(n-20) +57276494053174*a(n-21) -62381329826018*a(n-22) +61173951385090*a(n-23) -53931666825727*a(n-24) +42653735964962*a(n-25) -30177492913372*a(n-26) +19030922520904*a(n-27) -10649333267120*a(n-28) +5257938886816*a(n-29) -2274410125504*a(n-30) +854313663104*a(n-31) -275508944128*a(n-32) +75166819840*a(n-33) -17011807232*a(n-34) +3107903488*a(n-35) -440373248*a(n-36) +45408256*a(n-37) -3031040*a(n-38) +98304*a(n-39) for n>53
\end{quote}
As of the ``Last modified'' line on the live entry (Jul 23 2025 12:53:04, revision 6) this is still
recorded as empirical, and nothing on the entry records it as proved.

\section{What is being counted}

The array has $8$ cells across and entries $x(i,j)\in\{0,1,2\}$.

Two derived arrays are formed. The first reads the horizontal statistic across each
neighbouring pair of a row,
\[
b(i,j)\;=\;\bigl[\,a+b\,\bigr]_{a=x(i,j),\,b=x(i,j-1)}\qquad (j\ge1),
\]
and must not decrease DOWN a column: $b(i,j)\le b(i{+}1,j)$. The second reads the vertical
statistic down each neighbouring pair of a column,
\[
c(i,j)\;=\;\bigl[\,\lvert a-b\rvert\,\bigr]_{a=x(i,j),\,b=x(i{-}1,j)}\qquad (i\ge1),
\]
and must not decrease ALONG a row: $c(i,j)\le c(i,j{+}1)$. In words, the horizontal statistic
is the sum of a cell and its left neighbour, and the vertical one is
the absolute difference of a cell and the cell above it.

Nothing in the condition is invariant under renaming the letters --- each statistic either does
arithmetic on them or compares them by size --- so no reduction to equality patterns is
available, and the walk is run over the alphabet the entry names.

\section{A bounded window}

\begin{lemma}\label{lem:win}
Both conditions are conditions on TWO consecutive array rows and nothing more.
\end{lemma}

\begin{proof}
$b(i,j)\le b(i{+}1,j)$ names four cells, all in rows $i$ and $i+1$. $c(i,j)\le c(i,j{+}1)$
names four cells, all in rows $i-1$ and $i$. No comparison reaches further.
\end{proof}

So a single row is a state. Let $M$ be the matrix on the $3^{8}=6561$ rows with
$M[r,s]=1$ when the pair $(r,s)$ passes both tests and $0$ otherwise. An array of $R$ rows is a
walk of $R-1$ steps and every row may begin or end one, so with $\iota=\tau=\mathbf{1}$,
\[
a(n)\;=\;\iota^{\!\top}M^{\,j}\tau ,\qquad S=6561.
\]

The walk index $j$ and the entry's own $n$ differ by a constant, read off by matching the
published terms rather than assumed. In particular $a$ is $C$-finite.

\section{The proof}

\begin{lemma}\label{lem:crit}
Let $q(t)=t^{39}-\sum_i c_it^{39-i}$ be the characteristic polynomial of the
conjectured recurrence. The residual $a(n)-\sum_ic_ia(n-i)$ equals
$\iota^{\!\top}M^{\,j}q(M)\tau$ for the corresponding walk index $j$, so the recurrence
holds from some point on if and only if $\iota^{\!\top}M^{j}q(M)\tau=0$ for all large $j$,
and it is enough to examine $j<S$.
\end{lemma}

\begin{proof}
Factor $M^{j}$ out of $M^{j+39}-\sum_ic_iM^{j+39-i}$. For the bound, $M^{S}$
is an integer combination of $I,M,\dots,M^{S-1}$ by Cayley--Hamilton, so the numbers
$u_j=\iota^{\!\top}M^{j}q(M)\tau$ obey the monic recurrence given by the characteristic
polynomial of $M$; $S$ consecutive zeros force every later one, and the last nonzero $u_j$
pins the threshold exactly.
\end{proof}

\begin{theorem}
\[
a(n)\;=\;56\,a(n-1) - 1522\,a(n-2) + 26750\,a(n-3) - 341788\,a(n-4) + 3383870\,a(n-5) - 27019370\,a(n-6) + 178816330\,a(n-7) - 1000398245\,a(n-8) + 4801705370\,a(n-9) - 19999827232\,a(n-10) + 72938367572\,a(n-11) - 234575011664\,a(n-12) + 669091458700\,a(n-13) - 1700431069076\,a(n-14) + 3864487952276\,a(n-15) - 7876529587291\,a(n-16) + 14429295594092\,a(n-17) - 23796815540850\,a(n-18) + 35368238093958\,a(n-19) - 47398136192484\,a(n-20) + 57276494053174\,a(n-21) - 62381329826018\,a(n-22) + 61173951385090\,a(n-23) - 53931666825727\,a(n-24) + 42653735964962\,a(n-25) - 30177492913372\,a(n-26) + 19030922520904\,a(n-27) - 10649333267120\,a(n-28) + 5257938886816\,a(n-29) - 2274410125504\,a(n-30) + 854313663104\,a(n-31) - 275508944128\,a(n-32) + 75166819840\,a(n-33) - 17011807232\,a(n-34) + 3107903488\,a(n-35) - 440373248\,a(n-36) + 45408256\,a(n-37) - 3031040\,a(n-38) + 98304\,a(n-39)
\]
for every $n>53$.
\end{theorem}

\begin{proof}
The vector $w=q(M)\tau$ was formed by $39$ matrix--vector products in exact integer
arithmetic and $\iota^{\!\top}M^{j}w$ evaluated for $j=0,1,\dots$, again exactly; those
integers vanish from the index corresponding to $n=53+1$ onwards and the run continued
until the working vector was identically zero, which settles every larger $j$ at once.
\end{proof}

\section{Verification}

Four checks, each independent of the algebra above.

First, the model reproduces all $16$ terms the entry publishes, exactly, in integer
arithmetic.

Second, the count was recomputed for the smallest arrays straight from the definition: fill the
cells in row major order and test each comparison the moment both of its runs are complete,
in the orientation the entry's own name writes. No transfer matrix and no transposition are
involved, and the published terms came back.

Third, the conjectured recurrence was evaluated directly on the published terms, with no
matrices involved, and holds wherever the proved range and the data overlap.

Fourth, the annihilation test of Lemma~\ref{lem:crit} was carried out in exact integer
arithmetic on the whole vector rather than on a sampled prefix.

\begin{thebibliography}{9}
\bibitem{oeis} The OEIS Foundation, \emph{The On-Line Encyclopedia of Integer
Sequences}, \url{https://oeis.org}, 2026. Sequence A250941.
\bibitem{stanley} R.~P.~Stanley, \emph{Enumerative Combinatorics}, Volume 1, 2nd ed.,
Cambridge University Press, 2011. (Transfer-matrix method, Section 4.7.)
\bibitem{horn} R.~A.~Horn and C.~R.~Johnson, \emph{Matrix Analysis}, 2nd ed., Cambridge
University Press, 2013.
\end{thebibliography}

\end{document}
